\documentclass[11pt]{article}

\usepackage[utf8]{inputenc}
\usepackage{authblk}
\usepackage{amssymb}
\usepackage{geometry}
\geometry{margin=1in}

\title{A Fine Balance: Exploring the Interplay of Reproductive Strategies and Growth Dynamics in Trees}

\date{}

\begin{document}
\maketitle
You might have noticed trees in an area bearing way more fruit than usual every few years, and wondered what's going on. This is a common reproductive strategy called masting. One hypothesis is a growth-reproduction trade-off: trees put more resources into seeds in masting years, often triggered by a climatic event , but growth and seed production can also respond to climate independently, which makes teasing apart the trade-off tricky. For this project, We are trying to use long-term paired data on tree growth and seed production, from species with strong masting patterns, collected across stands at different elevations, to test this hypothesis with a hierarchical model. This is very much a work in progress, and I'd love feedback on the model. Come check out my poster!
\end{document}